\documentclass[11pt]{article}
\usepackage[utf8]{inputenc}
\usepackage{amsmath,amssymb}
\usepackage{hyperref}
\title{Robust Connectomics Brain Mapping under Distribution Shift}
\author{Anonymous}
\date{}
\begin{document}
\maketitle

\begin{abstract}
This paper studies Connectomics Brain Mapping. Grounded in a real, citable literature \cite{ref1,ref2,ref3,ref4,ref5,ref6,ref7,ref8},
we propose a focused method and report \textbf{illustrative} experiments.
All references below are real and verifiable (OpenAlex/arXiv); the reported
numbers are simulated to illustrate intended behaviour.
\end{abstract}

\section{Introduction}
Connectomics Brain Mapping is an active research area. This work targets a concrete gap identified from
the recent literature and contributes a method together with an evaluation protocol.

\section{Related Work (real literature)}
The following works, retrieved from public bibliographic APIs, frame our study:
\begin{itemize}
\item Gordon et al. (2017) — "Precision Functional Mapping of Individual Human Brains", Neuron (cited 1549).
\item Yeh et al. (2013) — "Deterministic Diffusion Fiber Tracking Improved by Quantitative Anisotropy", PLoS ONE (cited 1185).
\item Smith et al. (2013) — "Functional connectomics from resting-state fMRI", Trends in Cognitive Sciences (cited 1034).
\item Fornito et al. (2013) — "Graph analysis of the human connectome: Promise, progress, and pitfalls", NeuroImage (cited 803).
\item Fornito et al. (2012) — "Schizophrenia, neuroimaging and connectomics", NeuroImage (cited 748).
\item Deco et al. (2014) — "Great Expectations: Using Whole-Brain Computational Connectomics for Understanding Neuropsychiatric Disorders", Neuron (cited 467).
\end{itemize}

\section{Method}
We model distribution shift explicitly and optimise a worst-case objective over an uncertainty set of plausible test distributions. We instantiate this idea for Connectomics Brain Mapping and analyse its cost/quality trade-off.

\section{Experiments (Illustrative)}
\begin{center}
\begin{tabular}{lcc}
System & Primary Metric & Efficiency \\ \hline
Strong baseline & 0.812 & 1.00$\times$ \\
Ablation & 0.828 & 0.92$\times$ \\
\textbf{Ours} & \textbf{0.851} & \textbf{0.71$\times$}
\end{tabular}
\end{center}
\emph{Metrics simulated to illustrate the intended effect; not measured.}

\section{Conclusion}
We connected a real literature to a concrete method for Connectomics Brain Mapping. Future work will
validate the illustrative results with real experiments.

\begin{thebibliography}{9}
\bibitem{ref1} Evan M. Gordon, Timothy O. Laumann, Adrian W. Gilmore et al.. Precision Functional Mapping of Individual Human Brains. \emph{Neuron}, 2017. DOI:10.1016/j.neuron.2017.07.011
\bibitem{ref2} Fang‐Cheng Yeh, Timothy Verstynen, Yibao Wang et al.. Deterministic Diffusion Fiber Tracking Improved by Quantitative Anisotropy. \emph{PLoS ONE}, 2013. DOI:10.1371/journal.pone.0080713
\bibitem{ref3} Stephen M. Smith, Diego Vidaurre, Christian F. Beckmann et al.. Functional connectomics from resting-state fMRI. \emph{Trends in Cognitive Sciences}, 2013. DOI:10.1016/j.tics.2013.09.016
\bibitem{ref4} Alex Fornito, Andrew Zalesky, Michael Breakspear. Graph analysis of the human connectome: Promise, progress, and pitfalls. \emph{NeuroImage}, 2013. DOI:10.1016/j.neuroimage.2013.04.087
\bibitem{ref5} Alex Fornito, Andrew Zalesky, Christos Pantelis et al.. Schizophrenia, neuroimaging and connectomics. \emph{NeuroImage}, 2012. DOI:10.1016/j.neuroimage.2011.12.090
\bibitem{ref6} Gustavo Deco, Morten L. Kringelbach. Great Expectations: Using Whole-Brain Computational Connectomics for Understanding Neuropsychiatric Disorders. \emph{Neuron}, 2014. DOI:10.1016/j.neuron.2014.08.034
\bibitem{ref7} Hugues Duffau. Stimulation mapping of white matter tracts to study brain functional connectivity. \emph{Nature Reviews Neurology}, 2015. DOI:10.1038/nrneurol.2015.51
\bibitem{ref8} Mayuresh S. Korgaonkar, Alex Fornito, Leanne M. Williams et al.. Abnormal Structural Networks Characterize Major Depressive Disorder: A Connectome Analysis. \emph{Biological Psychiatry}, 2014. DOI:10.1016/j.biopsych.2014.02.018
\end{thebibliography}
\end{document}
